\IfFileExists{stacks-project.cls}{%
\documentclass{stacks-project}
}{%
\documentclass{amsart}
}

% For dealing with references we use the comment environment
\usepackage{verbatim}
\newenvironment{reference}{\comment}{\endcomment}
%\newenvironment{reference}{}{}
\newenvironment{slogan}{\comment}{\endcomment}
\newenvironment{history}{\comment}{\endcomment}

% For commutative diagrams we use Xy-pic
\usepackage[all]{xy}

% We use 2cell for 2-commutative diagrams.
\xyoption{2cell}
\UseAllTwocells

% We use multicol for the list of chapters between chapters
\usepackage{multicol}

% This is generally recommended for better output
\usepackage{lmodern}
\usepackage[T1]{fontenc}

% For cross-file-references
\usepackage{xr-hyper}

% Package for hypertext links:
\usepackage{hyperref}

% For any local file, say "hello.tex" you want to link to please
% use \externaldocument[hello-]{hello}
\externaldocument[introduction-]{introduction}
\externaldocument[conventions-]{conventions}
\externaldocument[sets-]{sets}
\externaldocument[categories-]{categories}
\externaldocument[topology-]{topology}
\externaldocument[sheaves-]{sheaves}
\externaldocument[sites-]{sites}
\externaldocument[stacks-]{stacks}
\externaldocument[fields-]{fields}
\externaldocument[algebra-]{algebra}
\externaldocument[brauer-]{brauer}
\externaldocument[homology-]{homology}
\externaldocument[derived-]{derived}
\externaldocument[simplicial-]{simplicial}
\externaldocument[more-algebra-]{more-algebra}
\externaldocument[smoothing-]{smoothing}
\externaldocument[modules-]{modules}
\externaldocument[sites-modules-]{sites-modules}
\externaldocument[injectives-]{injectives}
\externaldocument[cohomology-]{cohomology}
\externaldocument[sites-cohomology-]{sites-cohomology}
\externaldocument[dga-]{dga}
\externaldocument[dpa-]{dpa}
\externaldocument[sdga-]{sdga}
\externaldocument[hypercovering-]{hypercovering}
\externaldocument[schemes-]{schemes}
\externaldocument[constructions-]{constructions}
\externaldocument[properties-]{properties}
\externaldocument[morphisms-]{morphisms}
\externaldocument[coherent-]{coherent}
\externaldocument[divisors-]{divisors}
\externaldocument[limits-]{limits}
\externaldocument[varieties-]{varieties}
\externaldocument[topologies-]{topologies}
\externaldocument[descent-]{descent}
\externaldocument[perfect-]{perfect}
\externaldocument[more-morphisms-]{more-morphisms}
\externaldocument[flat-]{flat}
\externaldocument[groupoids-]{groupoids}
\externaldocument[more-groupoids-]{more-groupoids}
\externaldocument[etale-]{etale}
\externaldocument[chow-]{chow}
\externaldocument[intersection-]{intersection}
\externaldocument[pic-]{pic}
\externaldocument[weil-]{weil}
\externaldocument[adequate-]{adequate}
\externaldocument[dualizing-]{dualizing}
\externaldocument[duality-]{duality}
\externaldocument[discriminant-]{discriminant}
\externaldocument[derham-]{derham}
\externaldocument[local-cohomology-]{local-cohomology}
\externaldocument[algebraization-]{algebraization}
\externaldocument[curves-]{curves}
\externaldocument[resolve-]{resolve}
\externaldocument[models-]{models}
\externaldocument[functors-]{functors}
\externaldocument[equiv-]{equiv}
\externaldocument[pione-]{pione}
\externaldocument[etale-cohomology-]{etale-cohomology}
\externaldocument[proetale-]{proetale}
\externaldocument[relative-cycles-]{relative-cycles}
\externaldocument[more-etale-]{more-etale}
\externaldocument[trace-]{trace}
\externaldocument[crystalline-]{crystalline}
\externaldocument[spaces-]{spaces}
\externaldocument[spaces-properties-]{spaces-properties}
\externaldocument[spaces-morphisms-]{spaces-morphisms}
\externaldocument[decent-spaces-]{decent-spaces}
\externaldocument[spaces-cohomology-]{spaces-cohomology}
\externaldocument[spaces-limits-]{spaces-limits}
\externaldocument[spaces-divisors-]{spaces-divisors}
\externaldocument[spaces-over-fields-]{spaces-over-fields}
\externaldocument[spaces-topologies-]{spaces-topologies}
\externaldocument[spaces-descent-]{spaces-descent}
\externaldocument[spaces-perfect-]{spaces-perfect}
\externaldocument[spaces-more-morphisms-]{spaces-more-morphisms}
\externaldocument[spaces-flat-]{spaces-flat}
\externaldocument[spaces-groupoids-]{spaces-groupoids}
\externaldocument[spaces-more-groupoids-]{spaces-more-groupoids}
\externaldocument[bootstrap-]{bootstrap}
\externaldocument[spaces-pushouts-]{spaces-pushouts}
\externaldocument[spaces-chow-]{spaces-chow}
\externaldocument[groupoids-quotients-]{groupoids-quotients}
\externaldocument[spaces-more-cohomology-]{spaces-more-cohomology}
\externaldocument[spaces-simplicial-]{spaces-simplicial}
\externaldocument[spaces-duality-]{spaces-duality}
\externaldocument[formal-spaces-]{formal-spaces}
\externaldocument[restricted-]{restricted}
\externaldocument[spaces-resolve-]{spaces-resolve}
\externaldocument[formal-defos-]{formal-defos}
\externaldocument[defos-]{defos}
\externaldocument[cotangent-]{cotangent}
\externaldocument[examples-defos-]{examples-defos}
\externaldocument[algebraic-]{algebraic}
\externaldocument[examples-stacks-]{examples-stacks}
\externaldocument[stacks-sheaves-]{stacks-sheaves}
\externaldocument[criteria-]{criteria}
\externaldocument[artin-]{artin}
\externaldocument[quot-]{quot}
\externaldocument[stacks-properties-]{stacks-properties}
\externaldocument[stacks-morphisms-]{stacks-morphisms}
\externaldocument[stacks-limits-]{stacks-limits}
\externaldocument[stacks-cohomology-]{stacks-cohomology}
\externaldocument[stacks-perfect-]{stacks-perfect}
\externaldocument[stacks-introduction-]{stacks-introduction}
\externaldocument[stacks-more-morphisms-]{stacks-more-morphisms}
\externaldocument[stacks-geometry-]{stacks-geometry}
\externaldocument[moduli-]{moduli}
\externaldocument[moduli-curves-]{moduli-curves}
\externaldocument[examples-]{examples}
\externaldocument[exercises-]{exercises}
\externaldocument[guide-]{guide}
\externaldocument[desirables-]{desirables}
\externaldocument[coding-]{coding}
\externaldocument[obsolete-]{obsolete}
\externaldocument[fdl-]{fdl}
\externaldocument[index-]{index}

% Theorem environments.
%
\theoremstyle{plain}
\newtheorem{theorem}[subsection]{Theorem}
\newtheorem{proposition}[subsection]{Proposition}
\newtheorem{lemma}[subsection]{Lemma}

\theoremstyle{definition}
\newtheorem{definition}[subsection]{Definition}
\newtheorem{example}[subsection]{Example}
\newtheorem{exercise}[subsection]{Exercise}
\newtheorem{situation}[subsection]{Situation}

\theoremstyle{remark}
\newtheorem{remark}[subsection]{Remark}
\newtheorem{remarks}[subsection]{Remarks}

\numberwithin{equation}{subsection}

% Macros
%
\def\lim{\mathop{\mathrm{lim}}\nolimits}
\def\colim{\mathop{\mathrm{colim}}\nolimits}
\def\Spec{\mathop{\mathrm{Spec}}}
\def\Hom{\mathop{\mathrm{Hom}}\nolimits}
\def\Ext{\mathop{\mathrm{Ext}}\nolimits}
\def\SheafHom{\mathop{\mathcal{H}\!\mathit{om}}\nolimits}
\def\SheafExt{\mathop{\mathcal{E}\!\mathit{xt}}\nolimits}
\def\Sch{\mathit{Sch}}
\def\Mor{\mathop{\mathrm{Mor}}\nolimits}
\def\Ob{\mathop{\mathrm{Ob}}\nolimits}
\def\Sh{\mathop{\mathit{Sh}}\nolimits}
\def\NL{\mathop{N\!L}\nolimits}
\def\CH{\mathop{\mathrm{CH}}\nolimits}
\def\proetale{{pro\text{-}\acute{e}tale}}
\def\etale{{\acute{e}tale}}
\def\QCoh{\mathit{QCoh}}
\def\Ker{\mathop{\mathrm{Ker}}}
\def\Im{\mathop{\mathrm{Im}}}
\def\Coker{\mathop{\mathrm{Coker}}}
\def\Coim{\mathop{\mathrm{Coim}}}

% Boxtimes
%
\DeclareMathSymbol{\boxtimes}{\mathbin}{AMSa}{"02}

%
% Macros for moduli stacks/spaces
%
\def\QCohstack{\mathcal{QC}\!\mathit{oh}}
\def\Cohstack{\mathcal{C}\!\mathit{oh}}
\def\Spacesstack{\mathcal{S}\!\mathit{paces}}
\def\Quotfunctor{\mathrm{Quot}}
\def\Hilbfunctor{\mathrm{Hilb}}
\def\Curvesstack{\mathcal{C}\!\mathit{urves}}
\def\Polarizedstack{\mathcal{P}\!\mathit{olarized}}
\def\Complexesstack{\mathcal{C}\!\mathit{omplexes}}
% \Pic is the operator that assigns to X its picard group, usage \Pic(X)
% \Picardstack_{X/B} denotes the Picard stack of X over B
% \Picardfunctor_{X/B} denotes the Picard functor of X over B
\def\Pic{\mathop{\mathrm{Pic}}\nolimits}
\def\Picardstack{\mathcal{P}\!\mathit{ic}}
\def\Picardfunctor{\mathrm{Pic}}
\def\Deformationcategory{\mathcal{D}\!\mathit{ef}}
\usepackage{cleveref}
\crefname{originalreference}{original reference}{original references}
\Crefname{originalreference}{Original reference}{Original references}
\makeatletter
\newlabel{item-bounded-dimension}{{0DC4}{source}{}{sourcepointer-1}{}}
\newlabel{item-bounded-dimension@cref}{{[originalreference][0][]0DC4}{[1][0][]source}{}{}{}}
\makeatother
\crefname{originalreference}{original reference}{original references}
\Crefname{originalreference}{Original reference}{Original references}
\makeatletter
\newlabel{item-local}{{0DC3}{source}{}{sourcepointer-2}{}}
\newlabel{item-local@cref}{{[originalreference][0][]0DC3}{[1][0][]source}{}{}{}}
\makeatother
\crefname{originalreference}{original reference}{original references}
\Crefname{originalreference}{Original reference}{Original references}
\makeatletter
\newlabel{item-restriction}{{0DC2}{source}{}{sourcepointer-3}{}}
\newlabel{item-restriction@cref}{{[originalreference][0][]0DC2}{[1][0][]source}{}{}{}}
\makeatother
\crefname{originalreference}{original reference}{original references}
\Crefname{originalreference}{Original reference}{Original references}
\makeatletter
\newlabel{lemma-Serre-subcat-cartesian-modules}{{0D7J}{source}{}{sourcepointer-4}{}}
\newlabel{lemma-Serre-subcat-cartesian-modules@cref}{{[originalreference][0][]0D7J}{[1][0][]source}{}{}{}}
\makeatother
\crefname{originalreference}{original reference}{original references}
\Crefname{originalreference}{Original reference}{Original references}
\makeatletter
\newlabel{lemma-hypercovering-descent-modules}{{0D9Z}{source}{}{sourcepointer-5}{}}
\newlabel{lemma-hypercovering-descent-modules@cref}{{[originalreference][0][]0D9Z}{[1][0][]source}{}{}{}}
\makeatother
\crefname{originalreference}{original reference}{original references}
\Crefname{originalreference}{Original reference}{Original references}
\makeatletter
\newlabel{lemma-hypercovering-equivalence-bounded-modules}{{0DA2}{source}{}{sourcepointer-6}{}}
\newlabel{lemma-hypercovering-equivalence-bounded-modules@cref}{{[originalreference][0][]0DA2}{[1][0][]source}{}{}{}}
\makeatother
\crefname{originalreference}{original reference}{original references}
\Crefname{originalreference}{Original reference}{Original references}
\makeatletter
\newlabel{lemma-sanity-check-simplicial-semi-representable}{{0DC0}{source}{}{sourcepointer-7}{}}
\newlabel{lemma-sanity-check-simplicial-semi-representable@cref}{{[originalreference][0][]0DC0}{[1][0][]source}{}{}{}}
\makeatother
\crefname{originalreference}{original reference}{original references}
\Crefname{originalreference}{Original reference}{Original references}
\makeatletter
\newlabel{remark-augmentation-ringed}{{0D9W}{source}{}{sourcepointer-8}{}}
\newlabel{remark-augmentation-ringed@cref}{{[originalreference][0][]0D9W}{[1][0][]source}{}{}{}}
\makeatother
\crefname{originalsection}{original section}{original sections}
\Crefname{originalsection}{Original section}{Original sections}
\makeatletter
\newlabel{section-simplicial-semi-representable}{{0D8A}{source}{}{sourcepointer-9}{}}
\newlabel{section-simplicial-semi-representable@cref}{{[originalsection][0][]0D8A}{[1][0][]source}{}{}{}}
\makeatother
\crefname{originallemma}{original lemma}{original lemmas}
\Crefname{originallemma}{Original lemma}{Original lemmas}
\makeatletter
\newlabel{sites-cohomology-lemma-cohomology-modules-abelian-agree}{{03FD}{source}{}{sourcepointer-10}{}}
\newlabel{sites-cohomology-lemma-cohomology-modules-abelian-agree@cref}{{[originallemma][0][]03FD}{[1][0][]source}{}{}{}}
\makeatother
\crefname{originallemma}{original lemma}{original lemmas}
\Crefname{originallemma}{Original lemma}{Original lemmas}
\makeatletter
\newlabel{sites-cohomology-lemma-equivalence-unbounded-one}{{0D7V}{source}{}{sourcepointer-11}{}}
\newlabel{sites-cohomology-lemma-equivalence-unbounded-one@cref}{{[originallemma][0][]0D7V}{[1][0][]source}{}{}{}}
\makeatother
\title{Derived-category equivalence for cartesian modules over a hypercovering: La\textasciicircum{}* with quasi-inverse Ra\_*}
\author{The Stacks Project Authors}
\date{stacks-project:spaces-simplicial.tex:local-snapshot-20261012; GNU FDL 1.2 or later}
\begin{document}
\maketitle
\noindent\textit{Editorial note.} All delivered mathematical text is the author's own, copied from the pinned Stacks Project chapter file sources/stacks-spaces-simplicial/pinned-original-source.tex; no statement or proof step was rewritten, reordered or removed.\par\medskip
\noindent\textit{Editorial note.} The counted proof is the author's complete argument: the parenthetical directing the reader to the modules-versus-abelian-sheaves cohomology comparison (source lines 5520-5521) is a supplementary cross-reference, not an omission; the covering argument for the essential image is written out in full.\par\medskip
\noindent\textit{Editorial note.} The label lemma-hypercovering-equivalence-modules names the counted statement, and the proof environment immediately following it in the source is delivered verbatim as its proof.\par\medskip
\noindent\textit{Editorial note.} The counted statement is carried by the author's own lemma environment in the chapter source.\par\medskip
\noindent\textit{Allowed standard prerequisites (editorial).}\begin{itemize}
\item Cartesian modules form a weak Serre subcategory. Same chapter, \textbackslash{}label\{lemma-Serre-subcat-cartesian-modules\}. Lines 5457-5460: the base weak Serre fact, its hypotheses holding by the discussion in Remark remark-augmentation-ringed.
\item The bounded modules equivalence for hypercoverings. Same chapter, \textbackslash{}label\{lemma-hypercovering-equivalence-bounded-modules\}. Line 5468: La\_*a\textasciicircum{}*F = F, the unit half carried over from the bounded version.
\item Module-level descent: cartesian modules are pulled back. Same chapter, \textbackslash{}label\{lemma-hypercovering-descent-modules\}. Line 5474: every cartesian O-module G is a\textasciicircum{}*F for some F.
\item The weak Serre conditions on the A\_U. Same chapter, conditions (\textbackslash{}ref\{item-restriction\}), (\textbackslash{}ref\{item-local\}), (\textbackslash{}ref\{item-bounded-dimension\}). Lines 5490-5494 and 5510-5519: restriction and localness close the essential-image argument; bounded-dimension feeds assumptions (4) and (5) of the criterion.
\item Cohomology on Sites, the unbounded equivalence criterion. Cohomology on Sites chapter, \textbackslash{}label\{sites-cohomology-lemma-equivalence-unbounded-one\}. Line 5497: the five-assumption criterion that upgrades the abelian-category equivalence to the derived equivalence.
\item Cohomology of modules agrees with cohomology of abelian sheaves. Cohomology on Sites chapter, \textbackslash{}label\{sites-cohomology-lemma-cohomology-modules-abelian-agree\}. Lines 5520-5521: supplementary cross-reference for the reader worried about the module vs abelian sheaf distinction.
\end{itemize}\medskip
\ifdefined\hypertarget\hypertarget{sourcepointer-1}{}\fi
\noindent\textit{Source reference.} Original reference 0DC4: sources/stacks-spaces-simplicial/pinned-original-source.tex, source lines 5317--5317. Prerequisite pointer invoked at proof line(s) 5437; it names a standard fact of the same chapter that the counted argument uses, and every step of the counted proof built on it is retained verbatim in the delivered body.\par\medskip
\ifdefined\hypertarget\hypertarget{sourcepointer-2}{}\fi
\noindent\textit{Source reference.} Original reference 0DC3: sources/stacks-spaces-simplicial/pinned-original-source.tex, source lines 5310--5310. Prerequisite pointer invoked at proof line(s) 5437; it names a standard fact of the same chapter that the counted argument uses, and every step of the counted proof built on it is retained verbatim in the delivered body.\par\medskip
\ifdefined\hypertarget\hypertarget{sourcepointer-3}{}\fi
\noindent\textit{Source reference.} Original reference 0DC2: sources/stacks-spaces-simplicial/pinned-original-source.tex, source lines 5305--5305. Prerequisite pointer invoked at proof line(s) 5436; it names a standard fact of the same chapter that the counted argument uses, and every step of the counted proof built on it is retained verbatim in the delivered body.\par\medskip
\ifdefined\hypertarget\hypertarget{sourcepointer-4}{}\fi
\noindent\textit{Source reference.} Original reference 0D7J: sources/stacks-spaces-simplicial/pinned-original-source.tex, source lines 2669--2669. Prerequisite pointer invoked at proof line(s) 5457; it names a standard fact of the same chapter that the counted argument uses, and every step of the counted proof built on it is retained verbatim in the delivered body.\par\medskip
\ifdefined\hypertarget\hypertarget{sourcepointer-5}{}\fi
\noindent\textit{Source reference.} Original reference 0D9Z: sources/stacks-spaces-simplicial/pinned-original-source.tex, source lines 4692--4692. Prerequisite pointer invoked at proof line(s) 5476; it names a standard fact of the same chapter that the counted argument uses, and every step of the counted proof built on it is retained verbatim in the delivered body.\par\medskip
\ifdefined\hypertarget\hypertarget{sourcepointer-6}{}\fi
\noindent\textit{Source reference.} Original reference 0DA2: sources/stacks-spaces-simplicial/pinned-original-source.tex, source lines 4754--4754. Prerequisite pointer invoked at proof line(s) 5470; it names a standard fact of the same chapter that the counted argument uses, and every step of the counted proof built on it is retained verbatim in the delivered body.\par\medskip
\ifdefined\hypertarget\hypertarget{sourcepointer-7}{}\fi
\noindent\textit{Source reference.} Original reference 0DC0: sources/stacks-spaces-simplicial/pinned-original-source.tex, source lines 4016--4016. Prerequisite pointer invoked at proof line(s) 5513; it names a standard fact of the same chapter that the counted argument uses, and every step of the counted proof built on it is retained verbatim in the delivered body.\par\medskip
\ifdefined\hypertarget\hypertarget{sourcepointer-8}{}\fi
\noindent\textit{Source reference.} Original reference 0D9W: sources/stacks-spaces-simplicial/pinned-original-source.tex, source lines 4073--4073. Prerequisite pointer invoked at proof line(s) 5441; it names a standard fact of the same chapter that the counted argument uses, and every step of the counted proof built on it is retained verbatim in the delivered body.\par\medskip
\ifdefined\hypertarget\hypertarget{sourcepointer-9}{}\fi
\noindent\textit{Source reference.} Original section 0D8A: sources/stacks-spaces-simplicial/pinned-original-source.tex, source lines 3902--3902. Prerequisite pointer invoked at proof line(s) 5507; it names a standard fact of the same chapter that the counted argument uses, and every step of the counted proof built on it is retained verbatim in the delivered body.\par\medskip
\ifdefined\hypertarget\hypertarget{sourcepointer-10}{}\fi
\noindent\textit{Source reference.} Original lemma 03FD: sources/stacks-chow/tags.txt, source lines 4488--4488. Prerequisite pointer invoked at proof line(s) 5520; it names a standard fact of a sibling chapter that the counted argument uses, and every step of the counted proof built on it is retained verbatim in the delivered body.\par\medskip
\ifdefined\hypertarget\hypertarget{sourcepointer-11}{}\fi
\noindent\textit{Source reference.} Original lemma 0D7V: sources/stacks-chow/tags.txt, source lines 16570--16570. Prerequisite pointer invoked at proof line(s) 5498; it names a standard fact of a sibling chapter that the counted argument uses, and every step of the counted proof built on it is retained verbatim in the delivered body.\par\medskip
\begin{lemma}
\label{lemma-hypercovering-equivalence-modules}
Let $(\mathcal{C}, \mathcal{O}_\mathcal{C})$ be a ringed site.
Assume given weak Serre subcategories
$\mathcal{A}_U \subset \textit{Mod}(\mathcal{O}_U)$
satisfying conditions (\ref{item-restriction}),
(\ref{item-local}), and (\ref{item-bounded-dimension}) above.
Assume $\mathcal{C}$ has equalizers and fibre products and
let $K$ be a hypercovering.
Let $((\mathcal{C}/K)_{total}, \mathcal{O})$ be as in
Remark \ref{remark-augmentation-ringed}.
Let $\mathcal{A}_{total} \subset \textit{Mod}(\mathcal{O})$
denote the weak Serre subcategory of cartesian $\mathcal{O}$-modules
$\mathcal{F}$ whose restriction $\mathcal{F}_n$ is in
$\mathcal{A}_{K_n}$ for all $n$ (as defined above).
Then the functor $La^*$ defines an equivalence
$$
D_\mathcal{A}(\mathcal{O}_\mathcal{C})
\longrightarrow
D_{\mathcal{A}_{total}}(\mathcal{O})
$$
with quasi-inverse $Ra_*$.
\end{lemma}

\begin{proof}
The cartesian $\mathcal{O}$-modules form a weak Serre subcategory by
Lemma \ref{lemma-Serre-subcat-cartesian-modules}
(the required hypotheses hold by the discussion in
Remark \ref{remark-augmentation-ringed}).
Since the restriction functor
$g_n^* : \textit{Mod}(\mathcal{O}) \to \textit{Mod}(\mathcal{O}_n)$
are exact, it follows that $\mathcal{A}_{total}$ is a weak Serre
subcategory.

\medskip\noindent
Let us show that $a^* : \mathcal{A} \to \mathcal{A}_{total}$
is an equivalence of categories with inverse given by $La_*$.
We already know that $La_*a^*\mathcal{F} = \mathcal{F}$ by the
bounded version
(Lemma \ref{lemma-hypercovering-equivalence-bounded-modules}).
It is clear that $a^*\mathcal{F}$ is in $\mathcal{A}_{total}$
for $\mathcal{F}$ in $\mathcal{A}$. Conversely, assume that
$\mathcal{G} \in \mathcal{A}_{total}$. Because $\mathcal{G}$
is cartesian we see that $\mathcal{G} = a^*\mathcal{F}$
for some $\mathcal{O}_\mathcal{C}$-module $\mathcal{F}$ by
Lemma \ref{lemma-hypercovering-descent-modules}.
We want to show that $\mathcal{F}$ is in $\mathcal{A}$.
Take $U \in \Ob(\mathcal{C})$. We have to show that the
restriction of $\mathcal{F}$ to $\mathcal{C}/U$ is in $\mathcal{A}_U$.
As usual, write $K_0 = \{U_{0, i}\}_{i \in I_0}$.
Since $K$ is a hypercovering, the map $\coprod_{i \in I_0} h_{U_{0, i}} \to *$
becomes surjective after sheafification. This implies there is
a covering $\{U_j \to U\}_{j \in J}$ and a map $\tau : J \to I_0$
and for each $j \in J$ a morphism $\varphi_j : U_j \to U_{0, \tau(j)}$.
Since $\mathcal{G}_0 = a_0^*\mathcal{F}$ we find
that the restriction of $\mathcal{F}$ to $\mathcal{C}/U_j$
is equal to the restriction of the $\tau(j)$th component of
$\mathcal{G}_0$ to $\mathcal{C}/U_j$ via the morphism
$\varphi_j : U_j \to U_{0, \tau(i)}$. Hence by
(\ref{item-restriction}) we find that $\mathcal{F}|_{\mathcal{C}/U_j}$
is in $\mathcal{A}_{U_j}$ and in turn by
(\ref{item-local}) we find that $\mathcal{F}|_{\mathcal{C}/U}$
is in $\mathcal{A}_U$.

\medskip\noindent
In particular the statement of the lemma makes sense.
The lemma now follows from Cohomology on Sites,
Lemma \ref{sites-cohomology-lemma-equivalence-unbounded-one}.
Assumption (1) is clear (see Remark \ref{remark-augmentation-ringed}).
Assumptions (2) and (3) we proved in the preceding paragraph.
Assumption (4) is immediate from (\ref{item-bounded-dimension}).
For assumption (5) let $\mathcal{B}_{total}$ be the set of
objects $U/U_{n, i}$ of the site $(\mathcal{C}/K)_{total}$
such that $U \in \mathcal{B}$ where $\mathcal{B}$ is as in
(\ref{item-bounded-dimension}). Here we use the description of
the site $(\mathcal{C}/K)_{total}$ given in
Section \ref{section-simplicial-semi-representable}.
Moreover, we set $\text{Cov}_{U/U_{n, i}}$ equal to $\text{Cov}_U$
and $d_{U/U_{n, i}}$ equal $d_U$ where $\text{Cov}_U$ and $d_U$
are given to us by (\ref{item-bounded-dimension}).
Then we claim that condition (5) holds with these choices.
This follows immediately from
Lemma \ref{lemma-sanity-check-simplicial-semi-representable}
and the fact that $\mathcal{F} \in \mathcal{A}_{total}$
implies $\mathcal{F}_n \in \mathcal{A}_{K_n}$ and hence
$\mathcal{F}_{n, i} \in \mathcal{A}_{U_{n, i}}$.
(The reader who worries about the difference between
cohomology of abelian sheaves versus cohomology
of sheaves of modules may consult Cohomology on Sites, Lemma
\ref{sites-cohomology-lemma-cohomology-modules-abelian-agree}.)
\end{proof}

\end{document}
